\documentclass[a4paper, 12pt]{article}

\usepackage{utils}

\renewcommand*{\today}{01 April 2026}

\begin{document}

\hotbox{Probabilités 1}{CM 9}{\today}

\section{Inégalités et théorème limite}

\subsection{Markov (Inégalité de Markov)}

\begin{lemme}{}{}
    Soit $X, Y$ deux variables aléatoires de même type, possèdant une espérance telle que
    $\forall \omega \in \Omega, \; Y(\omega) \leq X(\omega)$. Alors
    $$
    E(Y) \leq E(X)
    $$
\end{lemme}

\begin{demonstration}
    Soit $Z = X - Y$, $Z$ est une variable aléatoire positive ou nulle de même type que $X$ et $Y$.
    C'est à dire que $Z(\Omega) \subset \R_+$

    \stump{mieux écrire}{moche}
    \begin{description}
        \item[Cas discret] $E(Z) = E(X) - E(Y)$ d'une part, et d'autre part $E(Z) = \sum_{z \in Z(\Omega)} z \P(Z = z) \geq 0$. Donc $E(X) - E(Y) \geq 0$.
        \item[Cas continu] $E(Z) = E(X) - E(Y)$ d'une part, et d'autre part $E(Z) = \int_{z \in Z(\Omega)} z f_Z(z) \mathrm{d}z \geq 0$. Donc $E(X) - E(Y) \geq 0$. 
    \end{description}
\end{demonstration}

\begin{remarque}
    En pratique, on peut utiliser ce lemme pour tout type de variables aléatoires, mais
    il est plus facile de démontrer ce lemme pour des variables aléatoires de même type.
\end{remarque}

\begin{lemme}{}{}
    Soit $X$ une variable aléatoire et $A \subset \R$. Soit $Y = \mathds{1}_A(X)$, alors
    $$
    E(Y) = P(X \in A)
    $$
\end{lemme}

\begin{demonstration}
    On a $P(Y = 1) = P(X \in A) := p$ et $P(Y = 0) = P(X \notin A) = 1 - p$. 

    Donc $Y \hookrightarrow B(1, p)$ et $E(Y) = p = P(X \in A)$.
\end{demonstration}

\begin{theoreme}{}{}
    Soit $X$ une variable aléatoire positive ou nulle possédant une espérance.
    $$
    \forall M > 0, \quad \P(X \geq M) \leq \frac{E(X)}{M}.
    $$ 
\end{theoreme}

\begin{remarque}
    Ce théorème est indépendant de la loi de $X$, ce qui réduit sa précision.
\end{remarque}

\begin{demonstration}
    Soit $Y(\omega) = \begin{cases}
        M & \text{si } X(\omega) \geq M \\
        0 & \text{sinon}
    \end{cases}$

    Alors on a $Y = M \mathds{1}_{[M, +\infty)}(X)$, d'où (d'après le lemme précédent) $E(Y) = M P(X \geq M)$.

    On compare $X$ et $Y$, ainsi on a
    $$
    X - Y = \begin{cases}
        X - M & \text{si } X \geq M \\
        X & \text{sinon}
    \end{cases} \; \geq 0
    $$
    Ainsi $Y \leq X$, et d'après le premier lemme, $E(Y) \leq E(X)$, c'est à dire
    $$
    M P(X \geq M) \leq E(X) \iff P(X \geq M) \leq \frac{E(X)}{M}
    $$
\end{demonstration}

\subsection{Inégalité de Bienaymé-Tchebychev}

Le principe est de
mesurer la dispersion d'une variable aléatoire autour de sa moyenne, et caractériser cette dispersion en terme de variance.

\begin{theoreme}{Théorème de Bienaymé-Tchebychev}{}
    Soit $X$ une variable aléatoire possèdant une espérance $\mu$ et une variance $\sigma^2$. Alors
    $$
    \forall a > 0, \quad P(|X - \mu| \geq a) \leq \frac{\sigma^2}{a^2}
    $$
\end{theoreme}

\begin{demonstration}
    Fixons $a > 0$.
    Étudions l'événement
    $$
    \{|X - \mu| \geq a\} = \{(X - \mu)^2 \geq a^2\}
    $$
    Ainsi on a
    $$
    P(|X - \mu| \geq a) = P((X - \mu)^2 \geq a^2) \quad \text{et} \quad E((X - \mu)^2) = \sigma^2
    $$

    Et en utilisant l'inégalité de Markov, on trouve
    $$
    P(|X - \mu| \geq a) \leq \frac{E((X - \mu)^2)}{a^2} = \frac{\sigma^2}{a^2}
    $$
\end{demonstration}

\subsection{Loi faible des grands nombres}

\begin{theoreme}{}{}
    Soit $(X_n)_{n \geq 1}$ une suite de variables aléatoires, deux à deux indépendantes, de même loi ("i.i.d" pour "independent and identically distributed"), possédant une espérance $\mu$ et une variance $\sigma^2$.

    Soit la suite $(\overline{X}_n)_{n \geq 1}$ définie par $\overline{X}_n = \frac{1}{n} \sum_{k=1}^n X_k$.

    Alors
    $$
    \forall \varepsilon > 0, \quad P(|\overline{X}_n - \mu| \geq \varepsilon) \xrightarrow[n \to +\infty]{} 0
    $$
\end{theoreme}

\begin{demonstration}
    Fixons $\varepsilon > 0$.

    On a $E(\overline{X}_n) = \mu$ et $\mathrm{Var}(\overline{X}_n) = \frac{\sigma^2}{n}$.

    En utilisant l'inégalité de Bienaymé-Tchebychev, on trouve
    $$
    P(|\overline{X}_n - \mu| \geq \varepsilon) \leq \frac{\mathrm{Var}(\overline{X}_n)}{\varepsilon^2} = \frac{\sigma^2}{n \varepsilon^2}
    $$

    Ainsi $P(|\overline{X}_n - \mu| \geq \varepsilon) \xrightarrow[n \to +\infty]{} 0$.
\end{demonstration}

\end{document}